% part: intuitionistic-logic
% chapter: soundness-completeness
% section: truth-lemma

\documentclass[../../../include/open-logic-section]{subfiles}

\begin{document}

\olfileid[hi]{int}{sc}{tru}

\olsection{सत्यता लेम्मा}

निम्न लेम्मा विहित प्रतिरूप में संतुष्टि को इस बात से जोड़ता है कि कौन-से
!!{formula} अभाज्य समुच्चय~$\Delta$ के !!{element} हैं।

\begin{lem}\ollabel{lem:truth}
  यदि $\Delta$ अभाज्य है, तो $\mSat{M(\Delta)}{!A}[\sigma]$ तभी जब
  $\Delta(\sigma) \Proves !A$।
\end{lem}

\begin{proof}
  ~$!A$ पर आगमन द्वारा।
  \begin{enumerate}
    \item \indcase{!A}{\lfalse}{चूँकि $\Delta(\sigma)$ अभाज्य है, वह संगत
      है; अतः $\Delta(\sigma) \Proves/ \indfrm$। परिभाषा से
      $\mSat/{M(\Delta)}{\indfrm}[\sigma]$।}
      \item \indcase{!A}{p}{$\mSat{{}}{}$ की परिभाषा से
        $\mSat{M(\Delta)}{\indfrm}[\sigma]$ तभी जब $\sigma \in V(p)$;
        अर्थात $\Delta(\sigma) \Proves \indfrm$।}
      \item \indcase!{!A}{\lnot !B}{}
      \item \indcase{!A}{!B \land
        !C}{$\mSat{M(\Delta)}{\indfrm}[\sigma]$ तभी जब
        $\mSat{M(\Delta)}{!B}[\sigma]$ और
        $\mSat{M(\Delta)}{!C}[\sigma]$। आगमन-परिकल्पना से
        $\mSat{M(\Delta)}{!B}[\sigma]$ तभी जब
        $\Delta(\sigma) \Proves !B$, और~$!C$ के लिए भी इसी प्रकार। किंतु
        $\Delta(\sigma) \Proves !B$ और $\Delta(\sigma) \Proves !C$ तभी
        जब $\Delta(\sigma) \Proves \indfrm$।}
      \item \indcase{!A}{!B \lor
        !C}{$\mSat{M(\Delta)}{\indfrm}[\sigma]$ तभी जब
        $\mSat{M(\Delta)}{!B}[\sigma]$ या
        $\mSat{M(\Delta)}{!C}[\sigma]$। आगमन-परिकल्पना से यह तभी जब
        $\Delta(\sigma) \Proves !B$ या $\Delta(\sigma) \Proves !C$। अब
        दिखाना है कि यह तभी जब $\Delta(\sigma) \Proves \indfrm$। बाएँ से
        दाएँ दिशा स्पष्ट है। दाएँ से बाएँ दिशा इस तथ्य से मिलती है कि
        $\Delta(\sigma)$ अभाज्य है।}
      \item \indcase{!A}{!B \lif !C}{पहले बाएँ से दाएँ दिशा का
        प्रतिधनात्मक रूप लें। मान लें
        $\Delta(\sigma) \Proves/ !B \lif !C$। तब
        $\Delta(\sigma) \cup \{!B\} \Proves/ !C$ भी। चूँकि किसी~$n$ के
        लिए $\tuple{!B, !C}$, $\tuple{!B_n, !C_n}$ है, इसलिए
        $\Delta(\sigma.n) = (\Delta(\sigma) \cup \{!B\})^*$; और
        $\Delta(\sigma.n) \Proves !B$, पर
        $\Delta(\sigma.n) \Proves/ !C$। आगमन-परिकल्पना से
        $\mSat{M(\Delta)}{!B}[\sigma.n]$ और
        $\mSat/{M(\Delta)}{!C}[\sigma.n]$। चूँकि
        $R\sigma(\sigma.n)$, इसका अर्थ है कि
        $\mSat/{M(\Delta)}{\indfrm}[\sigma]$।

        अब मान लें $\Delta(\sigma) \Proves !B \lif !C$, और
        $R\sigma\sigma'$ हो। चूँकि
        $\Delta(\sigma) \subseteq \Delta(\sigma')$, इसलिए यदि
        $\Delta(\sigma') \Proves !B$, तो
        $\Delta(\sigma') \Proves !C$। दूसरे शब्दों में, $R\sigma\sigma'$
        वाले प्रत्येक $\sigma'$ के लिए या तो
        $\Delta(\sigma') \Proves/ !B$ अथवा
        $\Delta(\sigma') \Proves !C$। आगमन-परिकल्पना से इसका अर्थ है कि
        जब भी $R\sigma\sigma'$, तब या तो
        $\mSat/{M(\Delta)}{!B}[\sigma']$ अथवा
        $\mSat{M(\Delta)}{!C}[\sigma']$; अर्थात
        $\mSat{M(\Delta)}{\indfrm}[\sigma]$।}
  \end{enumerate}
\end{proof}

\end{document}
